\begin{abstract}

A simulação numérica é ferramenta central de análise e projeto em Engenharia Mecânica, mas seu uso no ensino de métodos numéricos esbarra em duas dificuldades recorrentes: as ferramentas mais capazes exigem licenciamento, infraestrutura dedicada e configuração de ambiente, e raramente expõem de forma transparente a formulação matemática e o código por trás de cada simulação. O problema atacado neste trabalho é, portanto, reduzir a barreira de acesso a um ambiente de simulação numérica didático, sem abrir mão do rigor técnico nem da transparência da implementação. Como resposta, apresenta-se o \textit{MinervaMesh}, uma plataforma que resolve Equações Diferenciais Parciais de interesse em Engenharia Mecânica pelo Método dos Elementos Finitos (MEF) e pelo Método das Diferenças Finitas (MDF). A contribuição articula dois aspectos complementares, subordinados a esse único objetivo: (i) a implementação e a verificação dos métodos numéricos, confrontadas prioritariamente com soluções analíticas conhecidas e, nos casos sem solução fechada, com \textit{benchmarks} consagrados da literatura; e (ii) a entrega dessas simulações integralmente no navegador, em arquitetura \textit{client-side} sobre \textit{WebAssembly} e \textit{PyScript}, sem instalação e com o código-fonte Python inspecionável em cada caso. Os resultados mostram que as equações discretizadas reproduzem corretamente as soluções de referência e que a plataforma é adequada ao ensino, à prototipagem e a demonstrações acadêmicas de mecânica computacional.

\end{abstract}
